\begin{RiemannExercise}\label{ex:lie-groups}
\begin{enumerate}[label=(\arabic*),ref=\arabic*,leftmargin=*]
\item\label{ex:products-of-lie-groups-and-tori}
\begin{enumerate}[label=(\alph*),leftmargin=*]
\item 给定两个 Lie 群 $G_1,G_2$, 证明 $G_1\times G_2$ (两个群的直积) 赋有乘积的标准可微结构后是 Lie 群.
\item 圆 $S^1$ 可以与绝对值为 $1$ 的复数集合等同. 证明 $S^1$ 是 Lie 群, 并由此得出 $n$ 维环面 $T^n\cong S^1\times\cdots\times S^1$ 也是 Lie 群.
\end{enumerate}
\item\label{ex:abelian-lie-groups-and-exponential}
\begin{enumerate}[label=(\alph*),leftmargin=*]
\item 证明 $(\mathbb{R}^n,+)$ 是 Lie 群, 确定其 Lie 代数并写出指数映射的表达式.
\item 证明若 $G$ 是 Abel Lie 群, 则对所有 $V,W\in\mathfrak g$, 有 $[V,W]=0$.
\end{enumerate}
\item\label{ex:affine-maps-lie-group} 可以将
\[
  H=\{(x,y)\in\mathbb{R}^2\mid y>0\}
\]
中的每个点与由 $h(t)=yt+x$ 给出的可逆仿射映射 $h:\mathbb{R}\to\mathbb{R}$ 等同. 这样的映射全体关于复合构成一个群, 因而这一等同在 $H$ 上诱导出一个群结构.
\begin{enumerate}[label=(\alph*),leftmargin=*]
\item 证明诱导出的群运算由
\[
  (x,y)\cdot(z,w)=(yz+x,yw),
\]
给出, 且 $H$ 赋予这一群运算后是 Lie 群.
\item 证明左平移映射 $L_{(x,y)}:H\to H$ 在点 $(z,w)\in H$ 处的导数, 在上述坐标下由矩阵
\[
  (\mathrm{d}L_{(x,y)})_{(z,w)}=\begin{pmatrix}y&0\\0&y\end{pmatrix}.
\]
表示. 由此得出, 由向量
\[
  V=\xi\frac{\partial}{\partial x}+\eta\frac{\partial}{\partial y}
   \in\mathfrak h\equiv T_{(0,1)}H\qquad(\xi,\eta\in\mathbb{R})
\]
确定的左不变向量场 $X^V\in\mathfrak{X}(H)$ 由下式给出:
\[
  X_{(x,y)}^V=\xi y\frac{\partial}{\partial x}+\eta y\frac{\partial}{\partial y}.
\]
\item 给定 $V,W\in\mathfrak h$, 计算 $[V,W]$.
\item 确定向量场 $X^V$ 的流, 并给出指数映射 $\exp:\mathfrak h\to H$ 的表达式.
\item 先证明 $H$ 是 $GL(2)$ 中由满足 $y>0$ 的矩阵
\[
  \begin{pmatrix}y&x\\0&1\end{pmatrix}
\]
组成的子群, 再验证你的结果.
\end{enumerate}
\item\label{ex:special-linear-two-diffeomorphism} 考虑群
\[
  SL(2)=\left\{\begin{pmatrix}a&b\\c&d\end{pmatrix}\mid ad-bc=1\right\}.
\]
我们已经知道它是 3 维流形. 令
\[
  a=p+q,\quad d=p-q,\quad b=r+s,\quad c=r-s,
\]
证明 $SL(2)$ 与 $S^1\times\mathbb{R}^2$ 微分同胚.
\item\label{ex:exponentials-of-noncommuting-matrices} 举出矩阵 $A,B\in\mathfrak{gl}(2)$ 的例子, 使得 $e^{A+B}\ne e^Ae^B$.
\item\label{ex:determinant-of-matrix-exponential} 对 $A\in\mathfrak{gl}(n)$, 考虑可微映射
\begin{align*}
  h:\mathbb{R}&\to\mathbb{R}\setminus\{0\}\\
  t&\mapsto\det e^{At}
\end{align*}
并证明:
\begin{enumerate}[label=(\alph*),leftmargin=*]
\item 该映射是 $(\mathbb{R},+)$ 到 $(\mathbb{R}\setminus\{0\},\cdot)$ 的群同态;
\item $h'(0)=\operatorname{tr}A$;
\item $\det(e^A)=e^{\operatorname{tr}A}$.
\end{enumerate}
\item\label{ex:non-surjective-special-linear-exponential}
\begin{enumerate}[label=(\alph*),leftmargin=*]
\item 若 $A\in\mathfrak{sl}(2)$, 证明存在 $\lambda\in\mathbb{R}\cup i\mathbb{R}$, 使得
\[
  e^A=\cosh\lambda\,I+\frac{\sinh\lambda}{\lambda}A.
\]
\item 证明 $\exp:\mathfrak{sl}(2)\to SL(2)$ 不是满射.
\end{enumerate}
\item\label{ex:free-action-on-punctured-plane} 考虑由下式定义的向量场 $X\in\mathfrak{X}(\mathbb{R}^2)$:
\[
  X=\sqrt{x^2+y^2}\,\frac{\partial}{\partial x}.
\]
\begin{enumerate}[label=(\alph*),leftmargin=*]
\item 证明 $X$ 的流定义了 $\mathbb{R}$ 在 $M=\mathbb{R}^2\setminus\{0\}$ 上的自由作用.
\item 描述商拓扑空间 $M/\mathbb{R}$. 上述作用是固有的吗?
\end{enumerate}
\item\label{ex:diagonal-circle-action-on-product-of-spheres} 设 $M=S^2\times S^2$, 考虑 $S^1$ 在 $M$ 上由下式给出的对角作用:
\[
  e^{i\theta}\cdot(u,v)=\bigl(e^{i\theta}\cdot u,e^{2i\theta}\cdot v\bigr),
\]
其中, 对 $u\in S^2\subset\mathbb{R}^3$ 及 $e^{i\beta}\in S^1$, $e^{i\beta}\cdot u$ 表示将 $u$ 绕 $z$ 轴旋转角度 $\beta$.
\begin{enumerate}[label=(\alph*),leftmargin=*]
\item 确定此作用的不动点.
\item 可能的非平凡稳定子有哪些?
\end{enumerate}
\item\label{ex:closed-lie-subgroup-action-is-proper} 设 $G$ 是 Lie 群, $H$ 是闭 Lie 子群, 即 $H$ 是 $G$ 的子群且也是 $G$ 的闭子流形. 证明 $H$ 在 $G$ 上由 $A(h,g)=h\cdot g$ 定义的作用既自由又固有.
\item\label{ex:grassmannian-as-a-quotient} (\emph{Grassmann 流形}) 考虑形如
\[
  \begin{pmatrix}A&0\\C&B\end{pmatrix}
\]
的可逆矩阵所成的集合 $H\subset GL(n)$, 其中 $A\in GL(k)$, $B\in GL(n-k)$, $C\in\mathcal M_{(n-k)\times k}$.
\begin{enumerate}[label=(\alph*),leftmargin=*]
\item 证明 $H$ 是 $GL(n)$ 的闭 Lie 子群. 因而 $H$ 在 $GL(n)$ 上的作用既自由又固有 [参见 Exercise~\ref{ex:lie-groups}(\ref{ex:closed-lie-subgroup-action-is-proper})].
\item 证明商流形
\[
  Gr(n,k):=GL(n)/H
\]
可以与 $\mathbb{R}^n$ 的 $k$ 维子空间所成的集合等同 (特别地, $Gr(n,1)$ 就是射影空间 $\mathbb{R}P^{n-1}$).
\item 流形 $Gr(n,k)$ 称为 $\mathbb{R}^n$ 中 $k$ 维平面的 \textbf{Grassmann 流形}. 它的维数是多少?
\end{enumerate}
\item\label{ex:differential-of-lie-group-homomorphism} 设 $G$ 与 $H$ 是 Lie 群, $F:G\to H$ 是 Lie 群同态. 证明:
\begin{enumerate}[label=(\alph*),leftmargin=*]
\item $(\mathrm{d}F)_e:\mathfrak g\to\mathfrak h$ 是 Lie 代数同态;
\item 若 $(\mathrm{d}F)_e$ 是同构, 则 $F$ 是局部微分同胚;
\item 若 $F$ 是满射的局部微分同胚, 则 $F$ 是覆盖映射.
\end{enumerate}
\item\label{ex:quaternions-and-covering-of-rotation-group}
\begin{enumerate}[label=(\alph*),leftmargin=*]
\item 证明 $\mathbb{R}\cdot SU(2)$ 是 $\mathcal M_{2\times 2}(\mathbb{C})$ 的 4 维实线性子空间, 在矩阵乘法下封闭, 并有一组基
\begin{align*}
  1&=\begin{pmatrix}1&0\\0&1\end{pmatrix},&
  i&=\begin{pmatrix}i&0\\0&-i\end{pmatrix},\\
  j&=\begin{pmatrix}0&1\\-1&0\end{pmatrix},&
  k&=\begin{pmatrix}0&i\\i&0\end{pmatrix},
\end{align*}
满足 $i^2=j^2=k^2=ijk=-1$. 因此, 这一空间可以与\textbf{四元数}等同 (参见第~\RGref{subsec:notes-topological-manifolds}~节). 证明 $SU(2)$ 可以与 Euclidean 范数等于 $1$ 的四元数等同, 因而与 $S^3$ 微分同胚.
\item 若 $n\in\mathbb{R}^3$ 是一个单位向量, 将其与实部为零的一个四元数等同, 证明
\[
  \exp\left(\frac{n\theta}{2}\right)
    =1\cos\left(\frac{\theta}{2}\right)+n\sin\left(\frac{\theta}{2}\right)
\]
也是单位四元数.
\item 再将 $\mathbb{R}^3$ 与实部为零的四元数等同, 证明映射
\begin{align*}
  \mathbb{R}^3&\to\mathbb{R}^3\\
  v&\mapsto\exp\left(\frac{n\theta}{2}\right)\cdot v\cdot
       \exp\left(-\frac{n\theta}{2}\right)
\end{align*}
是绕 $n$ 确定的轴旋转角度 $\theta$ 的旋转.
\item 证明存在满射同态 $F:SU(2)\to SO(3)$, 并借此得出 $SU(2)$ 是 $SO(3)$ 的万有覆盖.
\item $SO(3)$ 的基本群是什么?
\end{enumerate}
\end{enumerate}
\end{RiemannExercise}
